\subsection{非線形の更新が元に戻るとき}\label{節：非線形の周期性}
  \sectionref*{節：初項依存型}では定数解や周期的な解を初項から調べた.
  ここでは,二つの初期値から始まる非線形の更新を文字のまま追い,
  一般項を先に求めずに周期を証明する.

  \[
    \begin{aligned}
      a_1&=x>0,\qquad a_2=y>0,\\
      a_{n+2}&=\frac{1+a_{n+1}}{a_n}\quad(n\ge1)
    \end{aligned}
  \]
  とする.正の二項から次も正になるため,すべての項が定義される.
  $x=y=1$で試すと$1,1,2,3,2,1,1,\ldots$となる.
  次は特定の数値から離れ,任意の正の$x,y$で計算しよう.

  \bookstep{最初の二項に戻るまでを計算する}
  \BookMathLayout{\[
    a_3=\frac{1+y}{x},\qquad
    a_4=\frac{1+x+y}{xy},\qquad
    a_5=\frac{1+x}{y}
  \]}{\[
    \begin{aligned}
      a_3&=\frac{1+y}{x},\\
      a_4&=\frac{1+x+y}{xy},\\
      a_5&=\frac{1+x}{y}
    \end{aligned}
  \]}
  である.例えば$a_5$では
  \[
    a_5=\frac{xy+1+x+y}{y(1+y)}=\frac{1+x}{y}
  \]
  と因数分解した.さらに
  \[
    a_6=\frac{1+a_5}{a_4}=x,\qquad
    a_7=\frac{1+a_6}{a_5}=y
  \]
  となる.二項の組が元に戻り,その組から先の項が一意に決まるので,
  \[
    a_{n+5}=a_n\qquad(n\ge1)
  \]
  である.この証明は正の初期値全体に使える.

  \bookstep{周期と最小周期を区別する}
  $5$は周期だが,常に最小周期というわけではない.
  定数列になるのは$x=y=t$, $t^2=t+1$のときであり,
  正の条件から$t=\frac{1+\sqrt5}{2}$である.このとき最小周期は$1$になる.
  ほかの初期値では最小周期は$5$である.
  なぜなら,最小周期は周期$5$の約数だからである.
  この約数の性質は,周期を最小周期で割った余りも周期になることから従う.

  \begin{bookpoint}
    \textbf{何が戻れば,繰り返しを証明できるか}\par
    次の項が直前の二項で決まるなら,一項だけでなく二項の組が戻ることを示す.
    初項から周期的な列,途中から周期的になる列,符号だけが振動する列は区別する.
  \end{bookpoint}
  \sectionref*{節：ガウス型}の有限個の状態や,
  \sectionref*{節：虚数解型隣接3項間漸化式}の回転とは別の仕組みで周期が現れた.
  数値の例と初項の例外を\bookproblemref{practice-connections}{5}{接続演習問5}で確かめよう.
